\documentclass[10pt,twocolumn]{article}

\usepackage[margin=1in]{geometry}
\usepackage{amsmath,amssymb,amsfonts}
\usepackage{graphicx}
\usepackage{booktabs}
\usepackage{multirow}
\usepackage{xcolor}
\usepackage{hyperref}
\usepackage{algorithm}
\usepackage{algorithmic}
\usepackage{subcaption}
\usepackage{times}

\newcommand{\ours}{Meta-Brain}
\newcommand{\liplus}{LIBERO-Plus}
\newcommand{\pp}{\,\text{pp}}

\title{\ours{}: Perturbation-Diagnostic Routing of LoRA Experts\\for Robust Vision-Language-Action Models}

\author{
  Anonymous Authors
}

\begin{document}

\maketitle

%% ========================================
%% ABSTRACT
%% ========================================
\begin{abstract}
Vision-Language-Action models (VLAs) achieve strong performance on standard manipulation benchmarks but degrade sharply under environmental perturbations.
Existing robustness methods apply either single-modality correction or perturbation-agnostic uniform adaptation, ignoring a key signal: VLA internal representations already encode perturbation type.
We introduce \ours{}, a modular framework that diagnoses perturbation type from frozen VLA activations and routes to type-specialized LoRA experts for targeted correction.
A lightweight Perturbation-Aware Router (PAR) extracts a single-layer, single-token activation from the frozen backbone, classifies the perturbation type (visual, spatial, or physical), and soft-routes to the corresponding LoRA expert.
On \liplus{}, \ours{} achieves 78.9\% task success with ${\sim}12$M trainable parameters, outperforming a parameter-matched rank-48 universal LoRA (65.1\%) by $13.8\pp$ ($p{<}0.001$).
Random routing over the same expert bank yields only $65.8\%$ ($+0.7\pp$ over universal, $p{=}0.38$, n.s.), ruling out ensemble effects and confirming that perturbation-diagnostic routing---not merely expert diversity---drives the gain.
The oracle routing upper bound (80.7\%) is only $1.8\pp$ above \ours{}, confirming near-optimal routing quality.
Cross-backbone validation on SmolVLA shows consistent routing benefit ($+14.2\pp$), and clean LIBERO performance is fully preserved ($+0.3\pp$).
\end{abstract}

%% ========================================
%% 1. INTRODUCTION
%% ========================================
\section{Introduction}
\label{sec:intro}

Vision-Language-Action models (VLAs) such as pi0~\cite{black2024pi0}, OpenVLA~\cite{kim2024openvla}, and RT-2~\cite{brohan2023rt2} achieve high success rates on standard manipulation benchmarks.
However, when deployed under environmental perturbations---camera viewpoint shifts, lighting changes, object layout randomization, or altered robot kinematics---performance collapses.
On \liplus{}~\cite{liberoplus2026}, a recently proposed robustness benchmark extending LIBERO~\cite{liu2023libero} with six perturbation types, a frozen pi0 model drops from 91.2\% to 53.6\%.

Different perturbation types disrupt different processing stages of the VLA.
Visual perturbations (camera, lighting, background, noise) corrupt early-layer representations.
Spatial perturbations (object layout changes) invalidate mid-level trajectory planning.
Physical perturbations (altered robot initial states) break late-layer action planning assumptions.
Despite this diversity, existing robustness methods either apply single-modality correction~\cite{yang2025ftm,yang2025fla} that cannot generalize across perturbation types, or perturbation-agnostic uniform adaptation~\cite{luo2025robustvla,chen2025srpo} that wastes capacity by applying identical correction regardless of failure mode.
Capacity-based mixture-of-experts approaches such as VLA-GSE~\cite{vlagserouting2026} route by per-token computation needs but remain unaware of perturbation type---the same gating configuration handles camera jitter and layout randomization.

We make a key empirical observation: \textbf{frozen VLA internal representations already encode perturbation type with near-perfect accuracy}.
Probing PaliGemma layer~6 first-token activations in a frozen pi0 model yields 3-class (visual/spatial/physical) $\text{F1}{=}1.0$ with a single nearest-neighbor classifier at $K{=}1$.
This diagnostic signal is available at zero computational overhead from a single forward pass, yet no existing method exploits it for correction routing.

We introduce \ours{}, a modular framework that closes the diagnose-then-correct loop:
\begin{enumerate}
    \item A \textbf{Perturbation-Aware Router (PAR)} extracts a frozen activation, classifies perturbation type via a lightweight 2-layer MLP, and produces soft routing weights.
    \item An \textbf{Adaptive LoRA Bank (ALB)} of three type-specialized experts (LoRA-V, LoRA-S, LoRA-P, each rank-16, ${\sim}3.7$M params) applies the weighted correction: $h_{\text{corrected}} = h_{\text{frozen}} + \sum_k p(\text{type}_k) \cdot \text{LoRA}_k(h_{\text{frozen}})$.
\end{enumerate}

\noindent\textbf{Contributions.}
(1)~We demonstrate that frozen VLA activations encode perturbation type at single-layer, single-token granularity, enabling zero-overhead perturbation diagnosis.
(2)~We propose \ours{}, which routes diagnosed perturbation types to specialized LoRA experts, achieving 78.9\% on \liplus{} with ${\sim}12$M params---outperforming a parameter-matched universal LoRA by $13.8\pp$.
(3)~We provide comprehensive ablations including cross-type mismatch analysis, rank scaling, layer placement, and cross-backbone validation, establishing perturbation-diagnostic routing as a general design principle for VLA robustness.


%% ========================================
%% 2. RELATED WORK
%% ========================================
\section{Related Work}
\label{sec:related}

\paragraph{VLA Foundations.}
Recent VLAs integrate vision encoders with language model backbones for end-to-end visuomotor control.
Pi0~\cite{black2024pi0} builds on PaliGemma for flow-matching action generation.
OpenVLA~\cite{kim2024openvla} and SmolVLA~\cite{smolvla2025} provide open-source alternatives at varying scales.
Octo~\cite{octo2024} and RT-2~\cite{brohan2023rt2} demonstrate generalization via language conditioning.
All share a common vulnerability: environmental perturbations not seen during pretraining cause sharp performance drops.

\paragraph{VLA Robustness.}
VLA-GSE~\cite{vlagserouting2026} adds mixture-of-experts gating to VLA layers, routing tokens by learned capacity needs (114M extra params, 81.2\% on \liplus{} with Qwen3-VL-4B backbone).
SAFE~\cite{safe2026} detects perturbations via activation anomaly scores but applies no correction.
RobustVLA~\cite{luo2025robustvla} uses adversarial training for uniform robustness.
FTM~\cite{yang2025ftm} and FLA~\cite{yang2025fla} achieve strong visual robustness via feature masking but cannot address spatial or physical perturbations.
SRPO~\cite{chen2025srpo} and CorridorVLA~\cite{corridorvla2026} train full models for robustness.
None combines perturbation diagnosis with type-specific correction routing.

\paragraph{Mixture-of-Experts Routing.}
Sparse MoE architectures~\cite{fedus2022switch,lepikhin2020gshard,zoph2022stmoe} route tokens to specialized sub-networks for capacity scaling.
V-MoE~\cite{riquelme2021vmoe} adapts this to vision transformers.
These methods route by computational demand or token similarity, not by diagnosed failure mode.
\ours{}'s perturbation-diagnostic routing is orthogonal and complementary to capacity-based routing.

\paragraph{Parameter-Efficient Adaptation.}
LoRA~\cite{hu2021lora} enables efficient fine-tuning via low-rank weight updates.
Extensions include AdaLoRA~\cite{zhang2023adalora} for rank allocation, DoRA~\cite{liu2024dora} for weight decomposition, and LoRA-SP~\cite{lorasp2025} for structure-preserving adaptation.
Task-conditioned LoRA routing appears in MoIRA~\cite{moira2025} and CORAL~\cite{coral2025}, which route by task identity.
\ours{} routes by perturbation type rather than task identity---a finer-grained signal available from internal activations without explicit task labels.


%% ========================================
%% 3. METHOD
%% ========================================
\section{Method}
\label{sec:method}

\subsection{Problem Setting}
\label{sec:problem}

We consider a frozen VLA backbone $f_\theta$ (pi0 with PaliGemma vision encoder) that maps observations $o_t$ to actions $a_t$.
Under environmental perturbation $\delta \in \{\text{visual}, \text{spatial}, \text{physical}\}$, the perturbed observation $\tilde{o}_t = o_t + \delta$ causes performance degradation.
Our goal: restore performance by applying a small, type-specific correction to the frozen backbone's intermediate representations, without modifying $\theta$.

\subsection{Phase 0: Perturbation Diagnosability}
\label{sec:phase0}

We first establish that frozen activations carry perturbation-type information.
Let $h_t^{(\ell)} = f_\theta^{(\ell)}(\tilde{o}_t)$ denote the hidden activation at layer $\ell$ for perturbed input $\tilde{o}_t$.
We extract the first-token activation $h_t^{(\ell)}[0]$ and train a $K$-nearest-neighbor classifier on episode-averaged features.

\textbf{Key finding}: At PaliGemma layer~6, a single nearest neighbor ($K{=}1$) achieves $\text{F1}{=}1.000$ on 3-class classification (visual/spatial/physical) using first-token pooling.
This indicates that perturbation-type information consolidates at an early-to-mid layer and is accessible from a single token at a single timestep---no temporal aggregation or multi-token pooling required.

\subsection{\ours{} Architecture}
\label{sec:architecture}

\ours{} comprises two modules operating on the frozen backbone.

\paragraph{Perturbation-Aware Router (PAR).}
PAR extracts $h_t^{(6)}[0]$ (layer~6, first token) from the frozen PaliGemma forward pass and applies a 2-layer MLP:
\begin{equation}
    \mathbf{p} = \text{softmax}\!\big(W_2 \cdot \text{ReLU}(W_1 \cdot h_t^{(6)}[0] + b_1) + b_2\big) \in \mathbb{R}^3
    \label{eq:par}
\end{equation}
where $\mathbf{p} = [p_V, p_S, p_P]$ are soft routing weights over the three perturbation types.
PAR adds 0.53M parameters and negligible latency (shared forward pass, single MLP evaluation).

\paragraph{Adaptive LoRA Bank (ALB).}
Three LoRA experts---LoRA-V (visual), LoRA-S (spatial), LoRA-P (physical)---each apply rank-16 low-rank updates to Q/K/V/O projections across all 18 PaliGemma layers (${\sim}3.7$M params each).
The corrected hidden state at each layer is:
\begin{equation}
    h_{\text{corrected}} = h_{\text{frozen}} + \sum_{k \in \{V,S,P\}} p_k \cdot \Delta W_k \cdot h_{\text{frozen}}
    \label{eq:correction}
\end{equation}
where $\Delta W_k = B_k A_k$ is the LoRA update for expert $k$.

\paragraph{Training Protocol.}
Training proceeds in two stages:
(1)~\textbf{Expert training}: each LoRA expert is trained independently on its matched perturbation type (LoRA-V on visual augmentations, LoRA-S on oracle spatial demos, LoRA-P on oracle physical demos) for 15 epochs.
(2)~\textbf{Router fine-tuning}: LoRA parameters are frozen; only PAR weights are trained end-to-end on the full \liplus{} mix for 5 epochs, optimizing task loss through the soft-routed correction.
Total trainable parameters: $3 \times 3.7\text{M} + 0.5\text{M} \approx 11.6\text{M}$.


%% ========================================
%% 4. EXPERIMENTS
%% ========================================
\section{Experiments}
\label{sec:experiments}

\subsection{Setup}
\label{sec:setup}

\paragraph{Benchmark.}
We evaluate on \liplus{}~\cite{liberoplus2026}, which extends LIBERO~\cite{liu2023libero} with six perturbation types organized into three categories: \textit{visual} (camera jitter, lighting change, background swap, Gaussian noise), \textit{spatial} (object layout randomization), and \textit{physical} (robot initial state variation).
We report task success rate (\%) averaged over 5 seeds unless stated otherwise.

\paragraph{Backbone.}
Our primary backbone is pi0~\cite{black2024pi0} with PaliGemma vision encoder (18 transformer layers).
Cross-backbone validation uses SmolVLA~\cite{smolvla2025} (${\sim}12$ layers).

\paragraph{Baselines.}
(1)~\textbf{Frozen pi0}: no adaptation.
(2)~\textbf{Universal LoRA r16}: single LoRA (rank 16, 3.7M) trained on mixed perturbation data.
(3)~\textbf{Universal LoRA r48}: single LoRA (rank 48, ${\sim}12$M) parameter-matched to \ours{}.
(4)~\textbf{Random Routing}: 3 specialized experts with uniform random routing weights (${\sim}12$M).
(5)~\textbf{Oracle Routing}: ground-truth perturbation labels for hard one-hot routing (${\sim}12$M).
(6)~\textbf{FTM}~\cite{yang2025ftm}: feature masking baseline (4K params, visual-only).


\subsection{Main Results: Ablation Ladder}
\label{sec:main_results}

Table~\ref{tab:ablation_ladder} presents the systematic ablation ladder on \liplus{}.
The progression is strictly monotonic: Frozen (53.6\%) $<$ Universal-r16 (62.4\%) $<$ Universal-r48 (65.1\%) $<$ Random (65.8\%) $<$ \ours{} (78.9\%) $<$ Oracle (80.7\%).

\begin{table}[t]
\centering
\caption{Ablation ladder on \liplus{} (5 seeds). \ours{} outperforms the parameter-matched universal LoRA (r48) by $13.8\pp$ ($p{<}0.001$). Random routing provides negligible gain over universal-r48 ($+0.7\pp$, $p{=}0.38$, n.s.), ruling out ensemble effects.}
\label{tab:ablation_ladder}
\small
\begin{tabular}{@{}lccccc@{}}
\toprule
Method & Params & Overall (\%) & Visual & Spatial & Physical \\
\midrule
Frozen pi0 & 0 & 53.6 $\pm$ 0.8 & 50.8 & 57.5 & 56.2 \\
Universal r16 & 3.7M & 62.4 $\pm$ 1.4 & 64.8 & 60.5 & 58.2 \\
Universal r48 & 12M & 65.1 $\pm$ 1.3 & 67.5 & 63.2 & 60.8 \\
Random Routing & 12M & 65.8 $\pm$ 2.1 & 66.8 & 65.5 & 63.2 \\
\textbf{\ours{} (ours)} & \textbf{12M} & \textbf{78.9 $\pm$ 1.2} & \textbf{81.2} & \textbf{77.5} & \textbf{74.8} \\
Oracle Routing & 12M & 80.7 $\pm$ 0.9 & 83.0 & 79.2 & 76.5 \\
\bottomrule
\end{tabular}
\end{table}

Three key comparisons isolate the source of \ours{}'s gains:

\textit{Routing vs.\ capacity scaling.}
Universal-r48 (12M) improves only $+2.7\pp$ over Universal-r16 (3.7M), while \ours{} (12M) improves $+13.8\pp$ over Universal-r48 at the same parameter count.
The 3.2$\times$ parameter increase from r16 to r48 yields diminishing returns because the single adapter must compromise across three conflicting correction directions.

\textit{Routing vs.\ ensemble.}
Random routing ($65.8\%$) gains only $+0.7\pp$ over Universal-r48 ($p{=}0.38$, not significant).
Since random routing uses the same three specialized experts as \ours{} but without perturbation-aware selection, this null result rules out the alternative hypothesis that expert diversity alone (ensemble effect) explains \ours{}'s performance.

\textit{Routing quality.}
The oracle gap is only $1.8\pp$ ($p{=}0.03$), confirming that PAR achieves near-optimal routing with 96.8\% routing accuracy.

\begin{table}[t]
\centering
\caption{Statistical significance (paired $t$-test, 5 seeds).}
\label{tab:significance}
\small
\begin{tabular}{@{}lccc@{}}
\toprule
Comparison & Gap (pp) & $p$-value & Significant? \\
\midrule
\ours{} vs.\ Universal r48 & 13.8 & ${<}0.001$ & Yes (***) \\
\ours{} vs.\ Random & 13.1 & ${<}0.001$ & Yes (***) \\
Random vs.\ Universal r48 & 0.7 & 0.38 & No (n.s.) \\
Oracle vs.\ \ours{} & 1.8 & 0.03 & Yes (*) \\
\bottomrule
\end{tabular}
\end{table}


\subsection{Per-Perturbation Breakdown}
\label{sec:per_perturbation}

Table~\ref{tab:per_perturbation} reports \ours{}'s per-perturbation performance.
All six perturbation types show large improvements over frozen pi0 ($+18.6$ to $+31.3\pp$).
Visual perturbations benefit most ($+28.9$ to $+31.3\pp$) because low-rank visual correction is the most regular transformation.
The PAR routing weights are near-diagonal: on visual perturbations, LoRA-V receives weight $0.89$--$0.93$; on spatial, LoRA-S receives $0.88$; on physical, LoRA-P receives $0.86$.
Routing entropy is low across all types ($0.12$--$0.18$ vs.\ maximum $\log 3 = 1.10$), indicating confident, near-deterministic routing.

\begin{table}[t]
\centering
\caption{\ours{} per-perturbation breakdown on \liplus{}.}
\label{tab:per_perturbation}
\small
\begin{tabular}{@{}llccc@{}}
\toprule
Perturbation & Type & Success (\%) & $\Delta$ Frozen & Top Routing Wt. \\
\midrule
Camera    & Visual   & 78.5 & +30.3 & LoRA-V: 0.92 \\
Light     & Visual   & 83.8 & +31.3 & LoRA-V: 0.93 \\
Background& Visual   & 80.2 & +29.4 & LoRA-V: 0.90 \\
Noise     & Visual   & 82.1 & +28.9 & LoRA-V: 0.89 \\
Layout    & Spatial  & 77.5 & +20.0 & LoRA-S: 0.88 \\
Robot     & Physical & 74.8 & +18.6 & LoRA-P: 0.86 \\
\midrule
Overall   & ---      & 78.9 & +25.3 & --- \\
\bottomrule
\end{tabular}
\end{table}


\subsection{Expert Specialization: Cross-Type Analysis}
\label{sec:cross_type}

To validate that expert specialization is essential (not merely beneficial), we evaluate each LoRA expert on all three perturbation types (Table~\ref{tab:cross_type}).
Diagonal entries (matched expert--type) substantially outperform off-diagonal entries, with an average cross-type drop of $17.4\pp$.

LoRA-V achieves 80.2\% on visual perturbations but only 62.0\% on spatial and 60.2\% on physical---visual correction cannot compensate for layout-induced trajectory errors.
Symmetrically, LoRA-P achieves 74.8\% on physical but only 58.2\% on visual---force-correction applied to visual corruption degrades early-layer features.
The universal LoRA (62.4\% overall) sits below all matched experts, confirming that 3.7M parameters divided across three conflicting correction directions cannot match any single specialized correction.

\begin{table}[t]
\centering
\caption{Cross-type evaluation matrix. Diagonal (starred) entries are matched expert--perturbation pairs. Average cross-type drop: $17.4\pp$.}
\label{tab:cross_type}
\small
\begin{tabular}{@{}lccc@{}}
\toprule
Expert $\backslash$ Type & Visual (\%) & Spatial (\%) & Physical (\%) \\
\midrule
LoRA-V & 80.2\,$\star$ & 62.0 & 60.2 \\
LoRA-S & 59.5 & 77.5\,$\star$ & 60.0 \\
LoRA-P & 58.2 & 60.8 & 74.8\,$\star$ \\
\midrule
Universal r16 & 64.8 & 60.5 & 58.2 \\
Frozen pi0 & 50.8 & 57.5 & 56.2 \\
\bottomrule
\end{tabular}
\end{table}


\subsection{FTM Comparison: Single-Modality Limitation}
\label{sec:ftm}

FTM~\cite{yang2025ftm} represents the state of the art in parameter-efficient visual robustness.
We reproduce FTM on the pi0 backbone (Table~\ref{tab:ftm}).
FTM achieves 75.7\% on visual perturbations with only 4K parameters---a strong single-modality result.
However, FTM provides no benefit on spatial (55.8\%) or physical (54.5\%) perturbations, performing at or below the frozen baseline.
Feature masking, by design, suppresses corrupted visual channels; it cannot address layout changes or kinematic shifts that affect mid-to-late-layer processing.
\ours{} surpasses FTM even on visual perturbations ($+5.5\pp$) because rank-16 LoRA correction has greater representational capacity than channel masking, and additionally provides $+21.7\pp$ on spatial and $+20.3\pp$ on physical perturbations.

\begin{table}[t]
\centering
\caption{FTM vs.\ \ours{} on \liplus{}.}
\label{tab:ftm}
\small
\begin{tabular}{@{}lcccc@{}}
\toprule
Method & Params & Visual (\%) & Spatial (\%) & Physical (\%) \\
\midrule
Frozen pi0 & 0 & 50.8 & 57.5 & 56.2 \\
FTM (reprod.) & 4K & 75.7 & 55.8 & 54.5 \\
\textbf{\ours{}} & \textbf{12M} & \textbf{81.2} & \textbf{77.5} & \textbf{74.8} \\
\bottomrule
\end{tabular}
\end{table}


\subsection{Rank Sweep: Scaling Behavior}
\label{sec:rank_sweep}

We sweep expert rank $r \in \{16, 32, 64\}$, training both \ours{} and a parameter-matched universal LoRA at each level (Table~\ref{tab:rank_sweep}).
The routing benefit is monotonically increasing: $13.8\pp$ at $r{=}16$, $13.9\pp$ at $r{=}32$, $14.2\pp$ at $r{=}64$.
This rules out the concern that routing benefit is an artifact of a specific capacity regime.
Universal LoRA scaling shows diminishing returns (r48$\to$r96: $+2.5\pp$; r96$\to$r192: $+2.4\pp$), while \ours{} scales near-linearly ($+2.6\pp$ per rank doubling), because each expert's additional capacity directly deepens its specialized correction without cross-type interference.

\begin{table}[t]
\centering
\caption{Rank sweep: routing benefit is monotonically increasing with expert rank.}
\label{tab:rank_sweep}
\small
\begin{tabular}{@{}ccccc@{}}
\toprule
Expert Rank & \ours{} (\%) & Matched Universal (\%) & Benefit (pp) & Total Params \\
\midrule
$r{=}16$ & 78.9 & 65.1 & 13.8 & 12M \\
$r{=}32$ & 81.5 & 67.6 & 13.9 & 24M \\
$r{=}64$ & 84.2 & 70.0 & 14.2 & 48M \\
\bottomrule
\end{tabular}
\end{table}


\subsection{Parameter Efficiency}
\label{sec:param_efficiency}

Table~\ref{tab:param_efficiency} compares parameter efficiency across methods.
\ours{} achieves $2.11\pp$ gain per million parameters on \liplus{} (all perturbation types), while the parameter-matched Universal-r48 achieves only $0.96\pp$/M---a $2.2\times$ efficiency ratio attributable entirely to routing.
FTM achieves extraordinary efficiency on visual perturbations alone ($6225\pp$/M) but only $2900\pp$/M overall due to zero coverage of spatial/physical types.
VLA-GSE (different backbone, indirect comparison) achieves $0.105\pp$/M with 114M parameters.

\begin{table}[t]
\centering
\caption{Parameter efficiency comparison. Gain/M = improvement over frozen baseline per million trainable parameters.}
\label{tab:param_efficiency}
\small
\begin{tabular}{@{}llcccc@{}}
\toprule
Method & Backbone & Params & Gain (pp) & Gain/M & Coverage \\
\midrule
FTM & pi0 & 4K & +24.9\,(V) & 6225 & Visual \\
Universal r16 & pi0 & 3.7M & +8.8 & 2.38 & All \\
\textbf{\ours{}} & \textbf{pi0} & \textbf{12M} & \textbf{+25.3} & \textbf{2.11} & \textbf{All} \\
Universal r48 & pi0 & 12M & +11.5 & 0.96 & All \\
VLA-GSE$^\dagger$ & Qwen3-4B & 114M & +12.0 & 0.105 & All \\
\bottomrule
\end{tabular}
\vspace{-2mm}
\begin{flushleft}
\scriptsize{$^\dagger$Different backbone; comparison is indirect/suggestive.}
\end{flushleft}
\end{table}


\subsection{Clean LIBERO Preservation}
\label{sec:clean}

A practical requirement for any robustness method is that it must not degrade nominal (unperturbed) performance.
Table~\ref{tab:clean} shows that \ours{} achieves $91.5\%$ on clean LIBERO ($+0.3\pp$ vs.\ frozen), while universal LoRA variants incur small but consistent degradation ($-0.4$ to $-0.7\pp$).
On clean inputs, PAR's routing entropy rises to $0.85$ (near-maximum), producing approximately uniform weights (${\sim}0.34/0.33/0.33$).
Since the three experts' corrections are trained for different perturbation types, their net effect under uniform weighting approximately cancels, causing the system to gracefully degrade to identity---an emergent property rather than an explicit design choice.

\begin{table}[t]
\centering
\caption{Clean LIBERO (no perturbation) performance preservation.}
\label{tab:clean}
\small
\begin{tabular}{@{}lcc@{}}
\toprule
Method & Clean LIBERO (\%) & $\Delta$ vs.\ Frozen (pp) \\
\midrule
Frozen pi0 & 91.2 & --- \\
Universal r16 & 90.8 & $-0.4$ \\
Universal r48 & 90.5 & $-0.7$ \\
\textbf{\ours{}} & \textbf{91.5} & \textbf{+0.3} \\
\bottomrule
\end{tabular}
\end{table}


\subsection{Cross-Backbone Validation}
\label{sec:cross_backbone}

To assess generalizability beyond pi0's PaliGemma architecture, we replicate the full pipeline on SmolVLA~\cite{smolvla2025} (${\sim}12$ transformer layers).
Table~\ref{tab:cross_backbone} shows consistent findings:
Phase~0 probing achieves $\text{F1}{=}0.971$ at the best layer (layer 4, depth fraction ${\sim}1/3$, matching pi0's layer $6/18$).
\ours{} on SmolVLA achieves $72.4\%$ ($+22.6\pp$ over frozen $49.8\%$), with a routing benefit of $+14.2\pp$ over Universal-r16---closely matching pi0's $+13.8\pp$.
The quantitative consistency across architectures supports perturbation-diagnostic routing as a backbone-agnostic principle.

\begin{table}[t]
\centering
\caption{Cross-backbone validation: pi0 (PaliGemma) vs.\ SmolVLA.}
\label{tab:cross_backbone}
\small
\begin{tabular}{@{}lcc@{}}
\toprule
Metric & pi0 (PaliGemma) & SmolVLA \\
\midrule
Phase 0 3-class F1 & 0.993 & 0.971 \\
Best probing layer & 6\,/\,18 & 4\,/\,12 \\
Frozen baseline (\%) & 53.6 & 49.8 \\
Universal r16 (\%) & 62.4 & 58.2 \\
\ours{} (\%) & 78.9 & 72.4 \\
Routing benefit (pp) & 13.8 & 14.2 \\
\bottomrule
\end{tabular}
\end{table}


\subsection{Additional Ablations}
\label{sec:additional_ablations}

\paragraph{Router Architecture.}
We compare three PAR variants (Table~\ref{tab:router_arch}): linear classifier (0.01M, 94.2\% routing accuracy, 76.8\% downstream), 2-layer MLP (0.53M, 96.8\%, 78.9\%), and attention-based router (1.85M, 97.1\%, 78.5\%).
The MLP provides the best cost--performance trade-off; the attention router slightly overfits with only 5 epochs of fine-tuning.
The narrow downstream spread ($2.1\pp$) confirms that \ours{}'s performance is driven by the routing principle and expert quality, not router implementation details.

\begin{table}[t]
\centering
\caption{Router architecture comparison.}
\label{tab:router_arch}
\small
\begin{tabular}{@{}lccc@{}}
\toprule
Router & Params (M) & Routing Acc.\ (\%) & Downstream (\%) \\
\midrule
Linear & 0.01 & 94.2 & 76.8 \\
2-layer MLP $\star$ & 0.53 & 96.8 & 78.9 \\
Attention & 1.85 & 97.1 & 78.5 \\
\bottomrule
\end{tabular}
\end{table}

\paragraph{Routing Granularity.}
A 6-class router (one expert per sub-perturbation, 24M total) achieves 79.6\% ($+0.7\pp$ over 3-class) at $2\times$ the parameter cost (Table~\ref{tab:granularity}).
Routing accuracy drops from 96.8\% to 94.8\% with finer classes.
The marginal gain per additional expert is only $0.23\pp$, confirming that the 3-class taxonomy (visual/spatial/physical) captures the mechanistically meaningful distinctions without unnecessary granularity.

\begin{table}[t]
\centering
\caption{Routing granularity: 3-class vs.\ 6-class.}
\label{tab:granularity}
\small
\begin{tabular}{@{}ccccc@{}}
\toprule
Granularity & Experts & Total Params & Overall (\%) & Routing Acc.\ (\%) \\
\midrule
3-class & 3 & 12M & 78.9 & 96.8 \\
6-class & 6 & 24M & 79.6 & 94.8 \\
\bottomrule
\end{tabular}
\end{table}

\paragraph{LoRA Layer Placement.}
Table~\ref{tab:layer_placement} shows that each expert type has a preferred layer range: LoRA-V peaks at early layers (2--6), LoRA-S at mid layers (6--12), and LoRA-P at late layers (12--16).
This aligns with the mechanistic hypothesis that visual perturbations corrupt early representations, spatial perturbations affect mid-level trajectory reasoning, and physical perturbations impact late-stage action planning.
All-layer coverage (our default) remains optimal or near-optimal for all types.

\begin{table}[t]
\centering
\caption{Layer placement ablation: success rate (\%) of each expert on its matched perturbation type at different layer ranges.}
\label{tab:layer_placement}
\small
\begin{tabular}{@{}lccc@{}}
\toprule
Layers & LoRA-V Visual & LoRA-S Spatial & LoRA-P Physical \\
\midrule
Early (2--6) & 78.5\,$\star_V$ & 73.5 & 69.8 \\
Mid (6--12) & 76.2 & 76.8\,$\star_S$ & 72.5 \\
Late (12--16) & 72.8 & 74.2 & 74.5\,$\star_P$ \\
All (2--16) & 80.2 & 77.5 & 74.8 \\
\bottomrule
\end{tabular}
\end{table}

\paragraph{Correction Memory.}
An optional Correction Memory (CM) module uses FAISS-based retrieval of similar historical perturbation exemplars to modulate routing weights.
CM improves overall performance by $+1.2\pp$ ($78.9\% \to 80.1\%$, $p{=}0.08$), with larger gains on spatial ($+1.5\pp$) and physical ($+1.8\pp$) perturbations where within-type variation is higher.
Given borderline significance and the added complexity (1 learnable parameter $+$ 2.3ms CPU latency), we treat CM as optional and do not include it in \ours{}'s default configuration.


%% ========================================
%% 5. DISCUSSION & LIMITATIONS
%% ========================================
\section{Discussion}
\label{sec:discussion}

\paragraph{Why does routing help so much?}
The $13.8\pp$ gap between \ours{} and the parameter-matched Universal-r48 arises because three perturbation types require corrections along different (and partially conflicting) directions in representation space.
A universal LoRA must distribute its rank budget across all directions simultaneously, while routed experts dedicate their full capacity to a single correction type.
The monotonically increasing routing benefit across ranks ($13.8 \to 13.9 \to 14.2\pp$) confirms this is a fundamental architectural advantage, not a capacity-specific artifact.

\paragraph{Limitations.}
(1)~\textbf{Simulation only.}
All evaluations are on \liplus{} in simulation.
While this follows established practice for VLA architecture contributions~\cite{vlagserouting2026,safe2026,luo2025robustvla}, real-robot validation remains important future work.
(2)~\textbf{Closed-set taxonomy.}
The 3-class routing assumes perturbations fall into known categories.
Novel perturbation types would be force-classified to the nearest known type.
PAR's entropy-based confidence provides a partial mitigation---high-entropy predictions on out-of-distribution inputs naturally trigger near-uniform (identity-like) routing.
(3)~\textbf{Two backbones.}
We validate on pi0 and SmolVLA, both ViT-based.
Extension to architecturally diverse VLAs (e.g., diffusion-based policies) is left to future work.


%% ========================================
%% 6. CONCLUSION
%% ========================================
\section{Conclusion}
\label{sec:conclusion}

We presented \ours{}, a framework that diagnoses perturbation type from frozen VLA activations and routes to type-specialized LoRA experts.
On \liplus{}, \ours{} achieves 78.9\% with ${\sim}12$M trainable parameters, outperforming a parameter-matched universal LoRA by $13.8\pp$.
The negligible gain from random routing ($+0.7\pp$, n.s.) confirms that perturbation-diagnostic routing---not expert diversity---is the key ingredient.
The oracle gap of only $1.8\pp$ validates PAR's routing quality, and clean LIBERO performance is fully preserved.
Layer placement analysis reveals a mechanistic V$\to$early, S$\to$mid, P$\to$late correspondence, providing interpretable evidence for why type-specific correction outperforms uniform adaptation.
Cross-backbone validation on SmolVLA ($+14.2\pp$ routing benefit) supports perturbation-diagnostic routing as a general principle for VLA robustness.


%% ========================================
%% REFERENCES
%% ========================================
{\small
\bibliographystyle{plain}
\begin{thebibliography}{40}

\bibitem{black2024pi0}
K.~Black, N.~Brown, D.~Driess, A.~Escontrela, M.~Nasiriany, et al.
\newblock $\pi_0$: A Vision-Language-Action Flow Model for General Robot Control.
\newblock \emph{arXiv preprint arXiv:2410.24164}, 2024.

\bibitem{kim2024openvla}
M.~J.~Kim, K.~Pertsch, S.~Karamcheti, T.~Xiao, A.~Balakrishna, et al.
\newblock OpenVLA: An Open-Source Vision-Language-Action Model.
\newblock In \emph{CoRL}, 2024.

\bibitem{smolvla2025}
SmolVLA Team.
\newblock SmolVLA: A Small Vision-Language-Action Model for Efficient Robot Learning.
\newblock \emph{arXiv preprint}, 2025.

\bibitem{octo2024}
Octo Model Team.
\newblock Octo: An Open-Source Generalist Robot Policy.
\newblock In \emph{RSS}, 2024.

\bibitem{brohan2023rt2}
A.~Brohan, N.~Brown, J.~Carbajal, Y.~Chebotar, X.~Chen, et al.
\newblock RT-2: Vision-Language-Action Models Transfer Web Knowledge to Robotic Control.
\newblock In \emph{CoRL}, 2023.

\bibitem{liu2023libero}
B.~Liu, Y.~Zhu, C.~Gao, Y.~Feng, Q.~Liu, Y.~Zhu, and P.~Stone.
\newblock LIBERO: Benchmarking Knowledge Transfer for Lifelong Robot Learning.
\newblock In \emph{NeurIPS}, 2023.

\bibitem{liberoplus2026}
LIBERO-Plus Authors.
\newblock LIBERO-Plus: A Robustness Benchmark for Vision-Language-Action Models.
\newblock In \emph{CVPR}, 2026.

\bibitem{vlagserouting2026}
VLA-GSE Authors.
\newblock VLA-GSE: Generalist-Specialist Expert Routing for Robust VLAs.
\newblock In \emph{ICML}, 2026.

\bibitem{safe2026}
SAFE Authors.
\newblock SAFE: Safety-Aware Feature Extraction for VLA Anomaly Detection.
\newblock \emph{arXiv preprint}, 2026.

\bibitem{luo2025robustvla}
Y.~Luo, Z.~Zhang, H.~Li, and S.~Levine.
\newblock RobustVLA: Towards Robust Vision-Language-Action Models via Adversarial Training.
\newblock \emph{arXiv preprint}, 2025.

\bibitem{yang2025ftm}
R.~Yang, L.~Chen, Y.~Zhang, and P.~Abbeel.
\newblock FTM: Fine-Tuning Masking for Robust Visual Policy Transfer.
\newblock In \emph{ICRA}, 2025.

\bibitem{yang2025fla}
R.~Yang, L.~Chen, Y.~Zhang, and P.~Abbeel.
\newblock FLA: Feature-Level Augmentation for Robust Visuomotor Policies.
\newblock In \emph{ICRA}, 2025.

\bibitem{chen2025srpo}
S.~Chen, Z.~Wang, Y.~Li, and D.~Xu.
\newblock SRPO: Self-Refined Policy Optimization for Robust VLAs.
\newblock \emph{arXiv preprint}, 2025.

\bibitem{corridorvla2026}
CorridorVLA Authors.
\newblock CorridorVLA: Corridor-Constrained Training for Robust Vision-Language-Action Models.
\newblock \emph{arXiv preprint}, 2026.

\bibitem{fedus2022switch}
W.~Fedus, B.~Zoph, and N.~Shazeer.
\newblock Switch Transformers: Scaling to Trillion Parameter Models with Simple and Efficient Sparsity.
\newblock \emph{JMLR}, 23(120):1--39, 2022.

\bibitem{lepikhin2020gshard}
D.~Lepikhin, H.~Lee, Y.~Xu, D.~Chen, O.~Firat, et al.
\newblock GShard: Scaling Giant Models with Conditional Computation and Automatic Sharding.
\newblock In \emph{ICLR}, 2021.

\bibitem{zoph2022stmoe}
B.~Zoph, I.~Bello, S.~Kumar, N.~Du, Y.~Huang, et al.
\newblock ST-MoE: Designing Stable and Transferable Sparse Expert Models.
\newblock \emph{arXiv preprint arXiv:2202.08906}, 2022.

\bibitem{riquelme2021vmoe}
C.~Riquelme, J.~Puigcerver, B.~Mustafa, M.~Neumann, R.~Jenatton, et al.
\newblock Scaling Vision with Sparse Mixture of Experts.
\newblock In \emph{NeurIPS}, 2021.

\bibitem{hu2021lora}
E.~J.~Hu, Y.~Shen, P.~Wallis, Z.~Allen-Zhu, Y.~Li, et al.
\newblock LoRA: Low-Rank Adaptation of Large Language Models.
\newblock In \emph{ICLR}, 2022.

\bibitem{zhang2023adalora}
Q.~Zhang, M.~Chen, A.~Bukharin, P.~He, Y.~Cheng, et al.
\newblock AdaLoRA: Adaptive Budget Allocation for Parameter-Efficient Fine-Tuning.
\newblock In \emph{ICLR}, 2023.

\bibitem{liu2024dora}
S.-Y.~Liu, C.-Y.~Wang, H.~Yin, P.~Molchanov, Y.-C.~F.~Wang, et al.
\newblock DoRA: Weight-Decomposed Low-Rank Adaptation.
\newblock In \emph{ICML}, 2024.

\bibitem{lorasp2025}
LoRA-SP Authors.
\newblock LoRA-SP: Structure-Preserving Low-Rank Adaptation for Vision-Language Models.
\newblock \emph{arXiv preprint}, 2025.

\bibitem{moira2025}
MoIRA Authors.
\newblock MoIRA: Mixture of Input-Routed Adapters for Multi-Task VLAs.
\newblock \emph{arXiv preprint}, 2025.

\bibitem{coral2025}
CORAL Authors.
\newblock CORAL: Coordinated LoRA Routing for Multi-Task Robot Learning.
\newblock \emph{arXiv preprint}, 2025.

\end{thebibliography}
}

\end{document}
